% =====================================================================
% PAPER A.  Photon-number parity in quantum spin ice.
% Draft.  Numbers marked \pending{} are NOT yet computed -- see README.md.
% =====================================================================
\documentclass[aps,prl,preprint,amsmath,amssymb,superscriptaddress,
longbibliography,floatfix]{revtex4-2}

\usepackage{bm}
\usepackage{graphicx}
\usepackage{xcolor}
\usepackage[colorlinks=true,citecolor=blue!55!black,
urlcolor=blue!55!black,linkcolor=blue!55!black]{hyperref}

\newcommand{\Jzz}{J_{zz}}
\newcommand{\Jpm}{J_{\pm}}
\newcommand{\Jpmpm}{J_{\pm\pm}}
\DeclareRobustCommand{\bq}{\bm q}
\DeclareRobustCommand{\br}{\bm r}
\newcommand{\PP}{\mathcal{P}}
% Every \pending{...} is a number or statement that has NOT been computed.
\setlength{\emergencystretch}{3em}
\newcommand{\pending}[1]{{\color{red}\textbf{[PENDING: #1]}}}

\begin{document}

\title{Photon-number parity in quantum spin ice:\\
exact selection rules for neutron, Raman and ultrasound spectroscopy}

\author{Zhengbang Zhou}
\email{zhengbang.zhou@mail.utoronto.ca}
\affiliation{Department of Physics, University of Toronto, Toronto,
Ontario M5S 1A7, Canada}

\author{Yong Baek Kim}
\email{ybkim@physics.utoronto.ca}
\affiliation{Department of Physics, University of Toronto, Toronto,
Ontario M5S 1A7, Canada}

\date{\today}

\begin{abstract}
The nearest-neighbour pseudospin Hamiltonian of quantum spin ice (QSI) is
invariant under a global $\pi$ rotation $\PP$ of the pseudospins about the
$\tilde x$ axis.  We show that on the ice manifold this operation is the
\emph{charge conjugation} of the emergent compact electrodynamics: it reverses
both the emergent electric and magnetic fields, so an $n$-photon state carries
parity $(-1)^n$ and photon-number parity is exactly conserved at zero magnetic
field.  Three consequences follow, none of which is perturbative or
model-detail dependent.  (i) Spectroscopic probes split into two closed
families: the emergent electric field $E=S^z$ and the linear magnetic field
$\sin B$ reach only odd-photon states, whereas every bond bilinear ---
$\cos B$, the Fleury--Loudon Raman operator, magnetostriction and the
ultrasound strain vertex --- reaches only even-photon states.  The one-photon
line and the two-photon continuum therefore appear in \emph{disjoint}
experiments.  (ii) At $\bq=0$ every single-site magnetic operator is a constant
on each electric-flux sector, so uniform magnetic probes are blind to the gauge
sector to all orders and only bond operators survive; this is why ultrasound,
not ESR or uniform susceptibility, is the natural $\bq\to0$ gauge probe.
(iii) A magnetic field is $\PP$-odd and switches the forbidden channels on at
order $h^2$, providing a parameter-free field protocol that distinguishes a
gauge excitation from any conventional mode.  We verify the selection rule on
the transport-masked FCC-32 spectrum, where every computed level is bright in
exactly one channel with a complement between $10^{-25}$ and $10^{-35}$, and we
show that the rule survives the pair-flip ($\Jpmpm$) term relevant to the Ce
pyrochlores.  In octupolar QSI the neutron dipole operator has no projection on
the gauge sector at all, so the even channel --- Raman and ultrasound --- is
the \emph{only} spectroscopic access to the emergent photon, and it sees the
two-photon continuum whose onset lies at twice the photon gap.
\end{abstract}

\maketitle

\section{Introduction}

The experimental case for quantum spin ice rests on identifying, in a measured
spectrum, an excitation that is the emergent photon of a compact $U(1)$ gauge
theory~\cite{Hermele2004,Gingras2014,Benton2012,RauGingras2019}.  In practice
this identification is made by exclusion: a low-lying feature that is not a
crystal-field level, not a spinon, and not a phonon is assigned to the gauge
sector.  Exclusion is weak evidence.  What is missing is a positive,
falsifiable statement of the form ``\emph{this} probe must see the mode and
\emph{that} probe must not.''

We provide such a statement.  It rests on a symmetry that is already present in
every nearest-neighbour pseudospin model used to describe the Ce and Yb
pyrochlores, and that has, to our knowledge, not been exploited
spectroscopically: the global $\pi$ rotation of the pseudospins about the
$\tilde x$ axis,
\begin{equation}
 \PP:\quad S^z_i\mapsto -S^z_i,\qquad S^{\pm}_i\mapsto S^{\mp}_i .
 \label{eq:PP}
\end{equation}
$\PP$ is one of the three $\mathbb{Z}_2$ elements of the Klein group of any
XYZ model with real couplings, and is therefore exact for the full
dipole--octupole Hamiltonian, not only for the XXZ limit.  Its gauge-theoretic
content is what matters here.  In the standard
dictionary~\cite{Hermele2004} the emergent electric field is $E_i = S^z_i$ and
the emergent vector potential is defined through $S^+_i = e^{iA_i}$, so
Eq.~(\ref{eq:PP}) reads
\begin{equation}
 \PP:\quad E\mapsto -E,\qquad A\mapsto -A,\qquad B\mapsto -B .
 \label{eq:C}
\end{equation}
This is charge conjugation.  Because the photon field is $C$-odd, an
$n$-photon state has $\PP$ eigenvalue $(-1)^n$: \textbf{photon-number parity is
an exactly conserved quantum number of quantum spin ice at zero field}.  The
statement is the lattice analogue of Furry's theorem in
QED~\cite{Furry1937}, and it is exact --- not a leading-order result of the
ring-exchange expansion, and not tied to any particular flux background.

The spectroscopic payoff is immediate and is the subject of this paper.  The
operators that experiments actually measure fall into two parity classes, and
the classes cannot talk to each other.  We state the classification in
Sec.~\ref{sec:rule}, verify it numerically in Sec.~\ref{sec:numerics} on
finite-cluster spectra from which the torus artifacts have been
removed~\cite{ZhouKimMask}, and translate it into a menu of experiments in
Sec.~\ref{sec:experiment}.  Section~\ref{sec:field} gives the field protocol
that turns the selection rule into a positive test rather than a null one.

\section{The selection rule}
\label{sec:rule}

\subsection{$\PP$ on the ice manifold}

Write the nearest-neighbour pseudospin Hamiltonian in the exchange-diagonal
local frame,
\begin{equation}
 H=\Jzz\sum_{\langle ij\rangle}S^z_iS^z_j
 -\Jpm\sum_{\langle ij\rangle}\left(S^+_iS^-_j+\mathrm{h.c.}\right)
 +\Jpmpm\sum_{\langle ij\rangle}\left(S^+_iS^+_j+\mathrm{h.c.}\right),
 \label{eq:model}
\end{equation}
with $\Jzz>0$.  Under Eq.~(\ref{eq:PP}) the Ising term is invariant, the
$\Jpm$ term maps to its Hermitian conjugate, and the $\Jpmpm$ pair-flip term
maps to its Hermitian conjugate as well.  Hence $[\PP,H]=0$ for the full XYZ
model.  In terms of the principal-axis components, $\PP$ acts as
$S^{\tilde x}\mapsto S^{\tilde x}$, $S^{\tilde y}\mapsto -S^{\tilde y}$,
$S^{\tilde z}\mapsto -S^{\tilde z}$, i.e.\ it is the $\pi$ rotation about
$\tilde x$.  The other two Klein elements are not independent on the ice
manifold: $R_{\tilde z}(\pi)=\exp[i\pi\sum_i(S^z_i+\tfrac12)]$ is a
$c$-number on any fixed-$S^z_{\rm tot}$ sector, so
$R_{\tilde y}(\pi)\equiv\PP$ there up to a phase.  \emph{One} nontrivial
$\mathbb{Z}_2$ acts on the gauge sector, and it is $\PP$.

Two structural remarks are needed before the rule can be used.

\emph{First}, $\PP$ reverses the total polarization
$\bm p=\sum_iS^z_i\br_i$ and therefore maps the electric-flux (transport)
sector $\mathcal T$ to $-\mathcal T$.  It is a symmetry \emph{within} a block
only for self-conjugate sectors, $\mathcal T=-\mathcal T$.  On FCC-32 the
ground state of the masked theory lies in the unique $\mathcal T=0$ sector, so
$\PP$ is a good quantum number of the ground state and of every state a
sector-preserving probe can reach.

\emph{Second}, both probes considered below \emph{are} sector preserving.  The
longitudinal probe $S^z_\mu(\bq)$ is diagonal in the Ising basis and cannot
change a configuration, let alone its transport label; the ring-exchange
operator around a \emph{contractible} hexagon returns the configuration to the
same sector by construction.  The selection rule therefore lives entirely
inside the $\mathcal T=0$ block, where $\PP$ is exact.  (Ring exchange around a
non-contractible six-cycle does change the sector; on FCC-32, $96$ of the $128$
six-cycles wrap the periodic box, and using them is exactly the torus artifact
that must be removed first~\cite{ZhouKimMask}.)

\subsection{Parity of the measurable operators}

Table~\ref{tab:parity} classifies the operators that appear in real
experiments.  The classification uses only Eq.~(\ref{eq:PP}) and is
independent of coupling strength, flux background and cluster.

\begin{table}[t]
\caption{Parity of gauge-sector operators under $\PP$, and the photon-number
change they can produce from a $\PP$-even ground state.  $W_{\mathcal C}
=S^+_1S^-_2\cdots S^-_L$ is the loop-flip (Wilson) operator, so
$W+W^\dagger=2\cos B$ and $i(W-W^\dagger)=-2\sin B$.}
\label{tab:parity}
\begin{ruledtabular}
\begin{tabular}{lccl}
operator & $\PP$ & $\Delta n$ & experimental channel\\ \hline
$E=S^z_\mu(\bq)$              & $-$ & odd  & neutron (dipolar QSI)\\
$\sin B \sim i(W-W^\dagger)$  & $-$ & odd  & --- (no linear-response probe)\\
$\cos B \sim W+W^\dagger$     & $+$ & even & Raman, THz two-photon\\
$S^+_iS^-_j+\mathrm{h.c.}$    & $+$ & even & Fleury--Loudon Raman~\cite{Fleury1968}\\
$S^z_iS^z_j$                  & $+$ & even & bond striction, ultrasound\\
$\varepsilon_{ab}\,\partial H/\partial\varepsilon_{ab}$ & $+$ & even & ultrasound, elastic constants\\
$h\sum_iS^z_i$                & $-$ & ---  & uniform field (breaks $\PP$)\\
\end{tabular}
\end{ruledtabular}
\end{table}

The content of the table is the statement
\begin{equation}
 \langle n|\,\mathcal O_+\,|0\rangle=0 \ \ \text{for odd } n,
 \qquad
 \langle n|\,\mathcal O_-\,|0\rangle=0 \ \ \text{for even } n,
 \label{eq:rule}
\end{equation}
where $\mathcal O_\pm$ is any $\PP$-even (odd) operator and $|n\rangle$ is an
$n$-photon state built on a $\PP$-even vacuum.  Three corollaries are worth
isolating.

\paragraph{(C1) The one-photon line and the two-photon continuum never appear
together.}  A neutron measurement of $S^{zz}(\bq,\omega)$ in a dipolar QSI
contains the single-photon pole and, at higher energy, the \emph{three}-photon
continuum --- but not the two-photon continuum, at any order in $\Jpm/\Jzz$.
Conversely a Raman or ultrasound spectrum contains the two-photon continuum
with \emph{no} one-photon line to sit on top of it.  The frequently invoked
worry that a broad two-magnon-like background contaminates the photon line, or
that an elastic one-photon peak swamps a weak continuum, is excluded by
symmetry.

\paragraph{(C2) $\langle\sin B\rangle=0$ and $\langle\cos B\rangle\neq0$ are
both consequences of the same rule.}  The measured hexagon flux
$\langle W+W^\dagger\rangle=\pm1/4$ of the $0$- and $\pi$-flux backgrounds is
$\PP$-even and may be finite; the conjugate quantity is forced to vanish.  A
nonzero $\langle\sin B\rangle$ in any calculation is a diagnostic of broken
$\PP$, i.e.\ of an applied field, a symmetry-breaking approximation, or a
finite-size artifact.

\paragraph{(C3) Photon decay $1\to2$ is exactly forbidden.}  The lowest
allowed photon-number-changing process is $1\to3$.  Together with the vanishing
of the polarization bubble below the two-spinon
threshold~\cite{Savary2012}, this makes the one-photon line intrinsically
sharp on symmetry grounds, not merely kinematically.  The leading photon
self-interaction is $2\to2$, whose strength is set by the emergent fine
structure constant~\cite{Pace2021}.

\subsection{The zero-momentum theorem}
\label{sec:qzero}

A second, independent superselection rule constrains the long-wavelength
limit.  On the ice manifold the four sublattice magnetizations
$m_\mu=\sum_{i\in\mu}S^z_i$ are \emph{constants} of each electric-flux sector.
We have verified on FCC-32 that the map from the $85$ transport sectors to the
$85$ magnetization tuples $(2m_0,2m_1,2m_2,2m_3)$ is a \emph{bijection}: the
transport label is not merely determined by the $\bq=0$ magnetizations, it
\emph{is} them (the same holds for the $25$ sectors of the $16$-site cubic
cluster).  The within-sector spread of every $2m_\mu$ is exactly zero, and the
commutator of each $S^z_\mu(\bq=0)$ with an arbitrary sector-block-diagonal
operator vanishes identically, against an unmasked control value of $55.9$.

Consequently every single-site magnetic operator at $\bq=0$ --- the uniform
magnetization, the staggered sublattice magnetizations, any ESR or uniform
susceptibility vertex --- is a $c$-number on the gauge sector and produces
\emph{zero} inelastic weight there, to all orders and at any coupling.  Only
operators that are bilinear in the spins on a bond, and therefore able to move
a configuration within a sector, couple to the gauge sector at $\bq\to0$.  This
is a symmetry-level explanation of why bond-striction and strain vertices,
rather than magnetic dipole vertices, are the correct $\bq\to0$ probes of the
photon~\cite{Patri2020,SimonPatriKim2022}.

\section{Numerical verification}
\label{sec:numerics}

The rule is exact, but it can only be \emph{exhibited} on a cluster whose gauge
sector is not contaminated.  A periodic pyrochlore cluster is a torus and
supports ice-preserving loops that close only through a periodic image; on
FCC-32 the shortest is a four-spin process at order $\Jpm^2$, one order
\emph{before} the physical hexagon at $|\Jpm|^3$~\cite{ZhouKimMask}.  Such a
process changes the transport sector, so the naive periodic spectrum has no
sharp $\PP$ classification at $\bq=0$ and no sharp flux sectors.  We therefore
work with the transport-masked Feshbach map of Ref.~\cite{ZhouKimMask}: virtual
spinon excursions are folded onto the $2970$-state FCC-32 ice manifold and all
matrix elements between different transport sectors are deleted.

We resolve every level of the masked block by cluster momentum with
$P_{\bq}=N_t^{-1}\sum_t\chi_{\bq}(t)^*T_t$ (projector exact to
$4\times10^{-13}$), and evaluate two channels: the longitudinal weight
$|\langle n|S^z_\mu(\bq)|0\rangle|^2$ and the magnetic weight
$|\langle n|\Phi(\bq)|0\rangle|^2$ with
$\Phi(\bq)=\sum_he^{-i\bq\cdot\br_h}(W_h+W_h^\dagger)$ summed over the $32$
\emph{contractible} hexagons only.

\begin{table}[t]
\caption{Masked FCC-32, virtual depth $d=1$.  Every level is bright in exactly
one channel.  $E$ is the longitudinal ($\PP$-odd) weight, $\cos B$ the
contractible ring-exchange ($\PP$-even) weight.  Energies in units of $\Jzz$,
measured from the ground state.}
\label{tab:levels}
\begin{ruledtabular}
\begin{tabular}{llll}
$\Jpm/\Jzz$ & $\omega$ \ (star) & $E$ weight & $\cos B$ weight\\ \hline
$-0.200$ & $0.015678$ \ ($L$) & $1.055$              & $9.9\times10^{-26}$\\
         & $0.025274$ \ ($X$) & $1.1\times10^{-27}$  & $1.645$\\
         & $0.033699$ \ ($L$) & $6.4\times10^{-26}$  & $0.201$\\
         & $0.037928$ \ ($X$) & $0.642$              & $9.2\times10^{-26}$\\ \hline
$+0.046$ & $0.002217$ \ ($L$) & $1.110$              & $1.6\times10^{-25}$\\
         & $0.003612$ \ ($X$) & $1.3\times10^{-26}$  & $2.891$\\
         & $0.004097$ \ ($L$) & $6.2\times10^{-27}$  & $0.560$\\
         & $0.004525$ \ ($X$) & $0.712$              & $9.2\times10^{-26}$\\
\end{tabular}
\end{ruledtabular}
\end{table}

Table~\ref{tab:levels} is the result.  The complement of the bright channel
lies between $10^{-25}$ and $10^{-35}$ throughout, on both sides of the sign of
$\Jpm$, i.e.\ in both the $\pi$-flux and the $0$-flux background.  The
selection is level by level, not channel by channel: every level the
longitudinal probe reaches has all four sublattice channels exactly equal, and
every level it misses is at machine zero in all four.  An earlier conjecture
that different irreducible representations of the little group carry the two
channels is thereby excluded --- the classification is the global parity, not a
spatial label.

Two further checks confirm that the rule is not an accident of the
construction.  Restricted to the $32$ contractible hexagons, the ring exchange
conserves the transport sector exactly (out-of-block weight $0.0$); including
the $96$ wrapping six-cycles instead moves $10$--$28\%$ of the flippable ice
weight out of the block, which is the artifact, correctly present.  And the
longitudinal $\bq=0$ response of Sec.~\ref{sec:qzero} is $1.42\times10^{-29}$
in the masked theory against $0.3126$ in the naive periodic theory at
$\Jpm/\Jzz=-0.20$ ($1.36\times10^{-30}$ against $0.4879$ at $-0.05$).

\pending{the $\PP$ classification has not yet been checked at $\Jpmpm\neq0$;
the proof in Sec.~\ref{sec:rule} covers it, but a numerical spot check at the
Ce$_2$Hf$_2$O$_7$ point $(\Jpm,\Jpmpm)/\Jzz=(-0.125,0.085)$ should be run.}

\subsection{The missing probe, and a two-photon bound state}
\label{sec:sinB}

Table~\ref{tab:levels} contains only one of the two odd operators.  Adding the
conjugate ring channel
\begin{equation}
 \Phi_-(\bq)=\sum_he^{-i\bq\cdot\br_h}\,i\left(W_h-W_h^\dagger\right)
 \ \sim\ \sin B ,
 \label{eq:sinB}
\end{equation}
which is $\PP$-\emph{odd}, completes the classification and makes a positive
photon assignment possible.  A single photon carries both an electric and a
magnetic amplitude, so it must be bright in $S^z$ \emph{and} in $\Phi_-$; a
two-photon state must be bright in $\cos B$ and dark in both.  The prediction
is therefore sharp:
\begin{equation}
 \text{levels bright in } S^z \ \Longleftrightarrow\ \text{bright in }\Phi_-
 \ \Longleftrightarrow\ \text{odd photon number.}
 \label{eq:assign}
\end{equation}
\pending{$\Phi_-$ has not been measured.  This is a few-hour run reusing
\texttt{campaign/run\_flux\_diagnostic.py}; see README.}

If Eq.~(\ref{eq:assign}) holds, Table~\ref{tab:levels} already contains a
striking hint.  On the fcc reciprocal lattice two $L$ points add to an $X$
point, $(1,1,1)+(1,-1,-1)=(2,0,0)$, so a non-interacting two-photon state at
$X$ built from the lowest $L$ excitation would sit at $2\omega(L)$.  Compare:
\begin{center}
\begin{tabular}{lccc}
$\Jpm/\Jzz$ & $2\omega(L)$ & $\omega_{\rm even}(X)$ & deficit\\ \hline
$-0.200$ & $0.031356$ & $0.025274$ & $19.4\%$\\
$+0.046$ & $0.004434$ & $0.003612$ & $18.5\%$\\
\end{tabular}
\end{center}
The lowest $\PP$-even level at $X$ lies below twice the lowest $\PP$-odd level
at $L$ by the same fraction, $\simeq19\%$, across a $4.3\times$ change of
coupling and across the flux-sector boundary.  If this is a two-photon bound
state, its binding energy is a direct nonperturbative measure of the emergent
photon--photon attraction --- i.e.\ of the emergent fine structure
constant~\cite{Pace2021} --- and it would appear in Raman and ultrasound as a
\emph{sharp line below the onset of the two-photon continuum}.

We flag two reasons this reading is a hypothesis and not yet a result.  First,
the ratio $\omega_{\rm odd}(X)/\omega(L)$ is $2.42$ and $2.04$ at the two
couplings, whereas a linear photon dispersion on this cluster would give
$|X|/|L|=2/\sqrt3=1.155$; the FCC-32 spectrum is far from the continuum limit.
Second, nothing here is converged in cluster size: between the $16$-site cubic
cluster and FCC-32 the lowest $X$ excitation moves by a factor $1.5$--$2.2$
and the $\Gamma$ excitation by $3.5$--$6.1$, and the $\Gamma/X$ ordering
inverts.  The selection rule (\ref{eq:rule}) is exact and size independent; the
energies it acts on are not.  The falsification test is stated in
Sec.~\ref{sec:tests}.

\section{Breaking the rule on purpose: the field protocol}
\label{sec:field}

A selection rule that only forbids is hard to use: an absent peak can always be
blamed on an intensity.  The remedy is that $\PP$ has a controlled breaking
parameter.  A uniform magnetic field enters as $-h\sum_i S^z_i$, which is
$\PP$-odd, so at finite $h$ the two families mix.  Standard perturbation theory
in the field gives, for a channel forbidden at $h=0$,
\begin{equation}
 I_{\rm forbidden}(h)= \left|\frac{h\,\langle n|M|n'\rangle
 \langle n'|\mathcal O|0\rangle}{E_n-E_{n'}}\right|^2 + O(h^4)
 \ \propto\ h^2 ,
 \label{eq:h2}
\end{equation}
with $M=\sum_iS^z_i$ the uniform magnetization.  The protocol is then:

\begin{enumerate}
\item At $h=0$, measure the Raman/ultrasound continuum (even) and the neutron
spectrum (odd).  The two must not share a feature.
\item Ramp $h$.  A gauge excitation acquires weight in its forbidden channel
with an intensity that grows as $h^2$ from exactly zero, with no threshold.
\item Any conventional excitation --- crystal-field level, phonon, magnon of an
ordered state, impurity mode --- has no reason to have zero weight at $h=0$ in
either channel, and no reason to turn on quadratically.
\end{enumerate}
The signature is the \emph{vanishing at $h=0$}, which is a property no
conventional mode shares.  Because $M$ is itself a sector label
(Sec.~\ref{sec:qzero}), the field also drives a staircase of transitions
between electric-flux sectors; on the $16$-site cluster the number of occupied
band sectors falls $25\to19\to16\to9$ as $h/\Jzz$ passes
$0.14$, $0.20$ and $0.28$ at $\Jpm/\Jzz=-0.05$, and the naive periodic
calculation delays this polarization crossover by a factor $2$--$3$ in $h$.
\pending{the $h^2$ onset of forbidden weight has not been computed; the field
term is not yet wired into the FCC-32 solvers (see README).}

\section{What each experiment sees}
\label{sec:experiment}

\subsection{Dipolar versus octupolar quantum spin ice}

For dipole--octupole doublets~\cite{HuangChenHermele2014} the neutron couples
only to the dipolar component of the pseudospin.  In the exchange-diagonal
frame the magnetic dipole is $m\propto\cos\theta\,\tilde\tau^z+
\sin\theta\,\tilde\tau^{x}$.  Within the ice manifold, $\tilde\tau^{x}$ flips a
spin and therefore leaves the manifold entirely --- it creates a spinon pair
and cannot contribute to the gauge sector.  The gauge-sector content of the
neutron vertex is thus $\cos\theta\,\tilde\tau^z=\cos\theta\,E$, a purely
$\PP$-odd operator.  Two limits follow:
\begin{itemize}
\item \textbf{Dipolar QSI} ($\cos\theta$ finite, ice built from a dipolar
component): neutrons see the one-photon line and the three-photon continuum,
and are blind to the two-photon continuum.
\item \textbf{Octupolar QSI} (ice built from the octupolar component): the
neutron vertex has \emph{no} projection on the gauge sector at all.  The photon
is neutron-dark for a reason stronger than a small form factor, and the even
channel is the only access.
\end{itemize}
The second case is the experimentally pressing one: the Ce pyrochlores
Ce$_2$Zr$_2$O$_7$ and Ce$_2$Hf$_2$O$_7$ are the leading QSI candidates
\cite{Gao2019,Gaudet2019,Bhardwaj2022,Smith2025} and the octupolar assignment
is under active discussion.  Our statement is then constructive rather than
pessimistic: \textbf{in octupolar QSI, Raman and ultrasound measure the
two-photon continuum, whose onset is at $2\omega_{\min}$ and therefore
determines the photon gap that neutrons cannot reach.}

\subsection{The menu}

\begin{table}[t]
\caption{Which probe reaches which part of the gauge sector.  ``--'' denotes an
exactly vanishing coupling at $h=0$, not a small one.}
\label{tab:menu}
\begin{ruledtabular}
\begin{tabular}{llll}
probe & vertex & parity & gauge sector, $h=0$\\ \hline
neutron, dipolar QSI      & $E(\bq)$        & $-$ & 1-photon, 3-photon\\
neutron, octupolar QSI    & spinon only     & --- & --\\
uniform susceptibility, ESR & $M$           & $-$ & -- (Sec.~\ref{sec:qzero})\\
Raman (Fleury--Loudon)    & $S_i\!\cdot\!S_j$ & $+$ & 2-photon continuum\\
THz two-photon absorption & $\cos B$        & $+$ & 2-photon\\
ultrasound / elastic const.& strain--bond   & $+$ & 2-photon\\
magnetostriction          & strain--bond    & $+$ & 2-photon\\
\end{tabular}
\end{ruledtabular}
\end{table}

Table~\ref{tab:menu} is the practical summary.  Three points deserve emphasis.
(i) The even-channel experiments do not compete with the one-photon line; they
measure a different object, and they measure it without an elastic background.
(ii) At $\bq\to0$, which is where ultrasound and Raman work, the magnetic
vertices are dead by the theorem of Sec.~\ref{sec:qzero}, so the bond/strain
vertex is not merely the largest coupling, it is the only one.  (iii) The
two-photon onset at $2\omega_{\min}$ is a \emph{lower} edge with a
characteristic density-of-states shape, and any sharp line below it is a bound
state of emergent light (Sec.~\ref{sec:sinB}).

\subsection{A warning about small-cluster calibrations}

Because the rule is exact, any calculation that violates it is diagnosing its
own artifacts.  Naive periodic exact diagonalization does violate it: the
winding processes mix transport sectors and convert the artifact into
observable $\bq=0$ spectral weight ($0.31$--$0.49$ of the sum rule on FCC-32,
against $\le10^{-29}$ after masking).  The same artifact inverts the measured
flux sign for every $\Jpm<0$, because the winding amplitude
$g_4=4\Jpm^2/\Jzz$ is even in $\Jpm$ while the physical hexagon amplitude
$g_6=12|\Jpm|^3/\Jzz^2$ is odd.  Published identifications calibrated on small
periodic clusters should be re-examined against these two diagnostics.

\section{Falsification tests}
\label{sec:tests}

The paper stands or falls on the following, in decreasing order of cost.

\begin{enumerate}
\item \textbf{$\Phi_-$ complementarity.}  Every level bright in $S^z$ must be
bright in $\Phi_-$ and dark in $\cos B$, and conversely.  Failure at machine
precision falsifies the identification $\PP=C$.
\item \textbf{Two-photon bound state.}  $2\omega_{\rm odd}(L)-\omega_{\rm
even}(X)$ must remain positive and approach a finite fraction as the cluster
grows, and the even level must be dark in $\Phi_-$.  If instead the deficit
tracks the finite-size drift of $\omega(X)$ it is not a bound state.
\item \textbf{$h^2$ onset.}  Forbidden-channel weight must vanish at $h=0$ and
rise quadratically, with the coefficient set by Eq.~(\ref{eq:h2}).
\item \textbf{$\Jpmpm$ robustness.}  The complementarity must survive the
pair-flip term at the Ce-relevant couplings, where $S^z_{\rm tot}$ is no longer
conserved.
\end{enumerate}

\section{Conclusion}

Quantum spin ice has an exact charge conjugation, inherited from the Klein
group of the XYZ pseudospin model, under which the emergent electric and
magnetic fields both change sign.  Photon-number parity is therefore an exact
quantum number at zero field, and the spectroscopic probes split into two
families that reach disjoint sets of excitations.  The one-photon line belongs
to neutrons; the two-photon continuum belongs to Raman, THz and ultrasound; a
uniform field breaks the rule at order $h^2$ and converts a null test into a
positive one.  In octupolar quantum spin ice the even channel is the only
access to the emergent photon at all.  These statements are exact, carry no
fitted parameter, and are independent of cluster size --- which matters,
because the energies they organize are not.

\begin{acknowledgments}
Computations used the transport-masked FCC-32 spectra of
Ref.~\cite{ZhouKimMask}.
\end{acknowledgments}

\bibliographystyle{apsrev4-2}
\bibliography{refs}

\end{document}
